% Internal research route, not a proof of the unrestricted conjecture.
\section{A coefficient-level route through growth deficits}
Let $E_w=1-z+(2z-1)F_w$. Its nonzero coefficients have interval
support $[\kappa(w),\operatorname{wt}(w)+1]$ for nonempty $w$.

\begin{proposition}[A sufficient coefficient bridge]
If $E_w$ is unimodal, then $F_w$ is unimodal.
\end{proposition}
\begin{proof}
Write $N=\operatorname{wt}w$, $q=z^NF_w(1/z)$, and
$e=z^{N+1}E_w(1/z)$. Then
\[
 e=(2-z)q-z^N+z^{N+1}.
\]
For $0\le r<N$ this gives $q_r=(e_r+q_{r-1})/2$, where
$q_{-1}=0$. Such geometric averaging preserves unimodality:
before a mode of $e$, the previous average is at most the new input,
so the averages increase. After that mode, once an average decreases,
it is still at least its current input, and the next input is no
larger; the averages keep decreasing. Thus $(q_0,\ldots,q_{N-1})$
is unimodal. Finally $q_{N-1}=q_N=1$, so appending $q_N$ preserves
unimodality. The empty host is immediate.
\end{proof}

\begin{proposition}[Appending a maximum preserves deficit unimodality]
Let $Z$ be a composition with $E_Z$ unimodal, and append
$t\ge\max Z$. Then $E_{Zt}$ is unimodal. In particular $E_Z$ is
unimodal for every weakly increasing $Z$.
\end{proposition}
\begin{proof}
Use the exact recurrence $E_{Zt}=S_tE_Z+z^{t+1}$.
The empty case is direct. Put $c=\operatorname{wt}Z$ and
$\kappa=\kappa(Z)$. The polynomial $W=S_tE_Z$ is unimodal.
If $c\ge t$, then $E_Z(t+1)\ge1$, since
$\kappa\le\max Z+1\le t+1\le c+1$. Moreover $E_Z(1)=0$.
Consequently $W_{t+2}-W_{t+1}=E_Z(t+1)-E_Z(1)\ge1$.
The added unit at $t+1$ cannot create a descent, and all other
differences keep the sign pattern of the unimodal $W$.
If $c<t$, the convolution $W$ is constant at its maximum throughout
$[c+2,t+\kappa]$. This interval contains $t+1$, because
$c+2\le t+1$ and $\kappa\ge2$. Adding a unit at one position of a
maximum plateau preserves unimodality. This proves the assertion.
\end{proof}

\paragraph{Two-arm identities.}
For $H=F_{XMY}$ and $A=F_X$, $B=F_Y$,
\[
 E_H=E_XD_Y+z^{M+1}B=E_YD_X+z^{M+1}A,
\]
and consequently
\[
 (1-z)E_H=E_XE_Y-z^{M+1}(2z-1)AB.
\]
These identities give a sufficient route in individual cases. The
unrestricted one-descent version of deficit unimodality is false,
as the explicit example below shows. Products of the displayed
factors need not be unimodal either.

\paragraph{An actual limitation, not an original-conjecture counterexample.}
Deficit unimodality is false for unrestricted hosts. The host
\[
 w=(1,1,1,1,1,1,3,1,2,1,2)
\]
has
\[
 E_w=(0,0,0,0,1,2,4,11,24,33,36,35,39,38,26,11,2)
\]
as its coefficient list, which decreases from $36$ to $35$ and
then increases to $39$. This host has two descents. Its rank
polynomial is still unimodal. The sufficiency proposition remains
valid, but any attempt to prove deficit unimodality for every host
would be false.

\paragraph{A one-descent obstruction.}
More decisively, take
\[
 w=(1,2,2,2,3,1,1,1,1,1,1,1).
\]
This host has exactly one descent and weight $17$. Its rank
polynomial and deficit have coefficient lists
\begin{align*}
 F_w&=(1,1,2,4,7,12,20,31,41,46,47,46,45,41,31,17,6,1),\\
 E_w&=(0,0,0,0,1,2,4,9,21,36,45,48,47,49,51,45,28,11,2).
\end{align*}
The deficit decreases from $48$ at rank $11$ to $47$ at rank
$12$, then increases to $49$ and $51$. The original coefficients
are unimodal, with their unique maximum $47$ at rank $10$.
Thus deficit unimodality cannot serve as a universal intermediate
theorem, even for peak three and an all-one second arm. The
sufficient bridge and the monotone-arm theorem proved above are
unaffected.

\begin{theorem}[No constant linear deficit gives a universal bridge]
For a real constant $c\ge2$, put
\[
 E^{(c)}_w=1-(c-1)z+(cz-1)F_w.
\]
For every such $c$, there is a one-descent host with unimodal rank
polynomial for which $E^{(c)}_w$ is not unimodal.
\end{theorem}
\begin{proof}
The weight-$17$ example above handles $c=2$. Suppose $c>2$.
Choose an integer $t\ge3$ with $t+4>2c$, and set
\[
 X=1^{t+3}t^t,\quad M=t+1,\quad
 a=\operatorname{wt}X=t^2+t+3,\quad
 \beta=a+M=t^2+2t+4,\quad Y=1^\beta.
\]
Let $A=F_X$ and $D=1+(S_M-1)A$. We claim that $D$ vanishes
below degree $M$ and has coefficients
\[
 (D_M,D_{M+1},D_{M+2})=(1,2,t+4).
\]
Here subscripts denote coefficient indices, with the parameter
$M$ held fixed.

To verify the claim, count the missing words of weights
$M,M+1,M+2$ in the ideal below $X$. Their numbers are respectively
$1,4,t+11$. At weight $M$ only the one-part word $(t+1)$ is
missing. At weight $M+1$, three words contain a part greater than
$t$, and the additional missing word is $(2,1^t)$.
At weight $M+2$, eight words contain a part greater than $t$.
Among the remaining words, a word fails to embed precisely when
the suffix beginning at its first non-one part has more than $t$
parts. There are two such words of length $t+2$, namely
$(2,1^{t+1})$ and $(1,2,1^t)$; there is one word $(3,1^t)$;
and there are $t$ words of length $t+1$ with two parts equal to
two, the first of them in position one. These are all possibilities.
The leading $t+3$ ones in $X$ absorb every needed initial run of
ones at these weights. Below $M$ every word embeds.

Using $(1-z)D=E^{(2)}_X-z^{M+1}A$ and the three missing counts,
the successive coefficients of $(1-z)D$ at $M,M+1,M+2$ are
$1,1,t+2$. This proves the claim about $D$.

Write $H=F_{XMY}$ and use $F_Y=[\beta+1]$. The maximum
decomposition gives the exact identity
\[
 (1-z)H=1+(S_M-z)A-z^{\beta+1}D.
\]
The middle term has degree $a+M=\beta$. Thus, at
$n_0=\beta+M+1$, the successive first differences of $H$ are
$-1,-2,-(t+4)$, while the preceding first difference is zero.
For indices here, the first difference of $E^{(c)}_w$ is
$c\,\Delta H_{n-1}-\Delta H_n$. At $n_0+1$ this equals
$2-c<0$, and at $n_0+2$ it equals $t+4-2c>0$. This is a
strict descent followed immediately by a strict ascent.
The original rank polynomial is unimodal by the all-one-arm
theorem. Thus the obstruction concerns the proposed auxiliary
deficit, not the one-descent rank-unimodality conjecture.
\end{proof}
